% 20260616_Assingment_Algebra_IIA
\documentclass[dvipdfmx]{jsreport}
\usepackage{soupreamble}
\title{代数学特論IIA 第2回レポート課題}
\author{M1 6012 川村聡一郎}
\date{}

\begin{document}
\maketitle
\begin{tcolorbox}
	$\wp(z)$が非線形微分方程式$(\wp'(z))^2= 4(\wp(z))^3- g_2(\wp(z))- g_3$の解になることを示す．\\
	ただし$g_2= 60\sum_{0\ne \omega\in \Lambda}\frac{1}{\omega^4},\ g_3= 140\sum_{0\ne \omega\in \Lambda}\frac{1}{\omega^6}$
\end{tcolorbox}

%\begin{souanswer} %解答はこっちに書く
    まず$\wp(z):= \frac{1}{z^2}+ \sum_{0\ne \omega\in \Lambda}\left(\frac{1}{(z- \omega)^2}- \frac{1}{\omega^2}\right)$を展開する．\\
    $\frac{1}{(z-\omega)^2} = \frac{1}{\omega^2}\left(1- \frac{z}{\omega}\right)^{-2}$であり，$(1- X)^{-2}= 1+ 2X+ 3X^2+ \cdots$と展開できるから，
\begin{equation*}
    \wp(z)= \frac{1}{z^2}+ \sum_{0\ne \omega\in \Lambda}\left( z\frac{2}{\omega^3}+ z^2\frac{3}{\omega^4}+ \cdots \right)
\end{equation*}
    Maclaurin展開の一意性と$\wp(z)= \wp(-z)$より
\begin{align*}
    \wp(z)&= \frac{1}{z^2}+ a_0+ a_1z+ a_2z^2+ a_3z^3+ \cdots\\
    \wp(-z)&= \frac{1}{z^2}+ a_0- a_1z+ a_2z^2- a_3z^3+ \cdots
\end{align*}
    したがって奇数番目の$a_k=0$．また$\wp(z)$の定義より$\lim_{z\to 0}\left(\wp(z)- \frac{1}{z^2}\right)= 0$であるから$a_0= 0$．\\
    $b_n= a_{2n}$とおくと$\wp(z)= \frac{1}{z^2}+ b_1z^2+ b_2z^4+ \cdots$とかけ，$b_n= (2n+ 1)\sum_{0\ne \omega\in \Lambda}\frac{1}{\omega^{2n+ 2}}$．\\
    各項を項別微分すると，
\begin{equation*}
    \wp'(z)= -\frac{2}{z^3}+ 2b_1z+ 4b_2z^3+ \cdots
\end{equation*}
    この$\wp'(z)$を2乗し，$\wp(z)$の3乗と比較すると，
\begin{equation*}
    (\wp'(z))^2- 4\wp(z)^3= -\frac{20b_1}{z^2}- 28b_2+ O(z^2)
\end{equation*}
    ここで問題文の定義より$g_2= 20b_1,\ g_3= 28b_2$となるから，
\begin{equation*}
    (\wp'(z))^2- 4\wp(z)^3= -\frac{g_2}{z^2}- g_3+ O(z^2)
\end{equation*}
    これを$F(z):= (\wp'(z))^2- 4\wp(z)^3+ g_2\wp(z)+ g_3$とおく．$\wp(z) = \frac{1}{z^2} + O(z^2)$より$g_2\wp(z) = \frac{g_2}{z^2} + O(z^2)$であるため，
\begin{equation*}
    F(z) = \left(-\frac{g_2}{z^2}- g_3+ O(z^2)\right) + \left(\frac{g_2}{z^2} + O(z^2)\right) + g_3 = O(z^2)
\end{equation*}
    となり，$F(z)$は$z=0$における特異点が解消され，$z=0$で正則かつ$F(0)=0$となる．\\
    また，$\wp(z), \wp'(z)$は周期格子$\Lambda$に関する楕円関数であるから，その多項式である$F(z)$もまた楕円関数である．$F(z)$の極となり得る点は$z \in \Lambda$のみであるが，二重周期性より$z=0$で正則であればすべての格子点でも正則となる．\\
    したがって$F(z)$は全複素平面で極を持たない二重周期関数（有界な整関数）となるため，Liouvilleの定理より定数関数である．
    $F(0) = 0$であることから，すべての$z \in \mathbb{C}$に対して$F(z) \equiv 0$となり，題意が示された．
%\end{souanswer}
\end{document}


% 間違ってるっぽいけど一応こっちも残しておく
\begin{souanswer} %解答はこっちに書く
	まず$\wp(z):= \frac{1}{z^2}+ \sum\left(\frac{1}{(z- \omega)^2- \frac{1}{\omega^2}$を展開する．\\
	$\frac{1}{1- X}= 1+ X+ X^2+ \cdots$と変形できるから，
\begin{equation*}
	\wp(z)= \frac{1}{z^2}+ \sum z\frac{2}{\omega^3}+ \sum z^2\frac{3}{\omega^4}+ \cdots
\end{equation*}
	Maclaurin展開の一意性と$wp(z)= \wp(-z)$より
\begin{align*}
	\wp(z)&= \frac{1}{z^2}+ a_0+ a_1z+ a_2z^2+ a_3z^3+ \cdots\\
	\wp(-z)&= \frac{1}{z^2}+ a_0- a_1z+ a_2z^2- a_3z^3+ \cdots\\
\end{align*}
	したがって奇数番目の$a_k=0$．また$\wp(z)$の定義より$\lim_{z\to 0}\left(\wp(z)- \frac{1}{z^2}\right)= 0$であるから$a_0= 0$\\
	$b_n= a_{2n}$とおくと$\wp(z)= \frac{1}{z^2}+ b_1z^2+ b_2z^4+ \cdots$とかけ，$b_n= (2n+ 1)\sum\frac{1}{\omega^{2n+ 2}}$．\\
	各項を項別微分すると，
\begin{equation*}
	\wp'(z)= -\frac{2}{x^3}+ 2b_1z+ 4b_2z^3+ \cdots+ 2nb_nz^{2n- 1}+ \cdots
\end{equation*}
	この$\wp'(z)$を2乗し，$\wp(z)$の3乗と比較すると，
\begin{equation*}
	(\wp'(z))^2- 4\wp^3(z)= -\frac{20}{z^2}b_1- 28b_2+ z^2(O(z^2))
\end{equation*}
	ここで$g_2= 20b\1, g_3= 28b\2$とおくと$(\wp'(z))^2- 4\wp^3(z)= -\frac{g_2}{z^2}- g_3+ z^2(O(z^2))$．またこれを$F(z):= (\wp'(z))^2- 4\wp^3(z)+ g_2\wp(z)+ g_3)$とする．$F(z)= 0$を示すためには$z= 0$であるすべての$z\in \C$に対して$O(z^2)\equiv 0$をいえばよい．\\
	$O(z^2)$は$O(z^2)= 12(b_1^2- b_3)+ c_2x^2+ c_4z^4+ \cdots$とかけ，係数がすべて0となればよいから$b_1^2- b_3= 0$となればよい．
\end{souanswer}
\end{document}
